\documentclass[a4paper,12pt]{article}
\usepackage[utf8]{inputenc}
\usepackage[T2A]{fontenc}
\usepackage[russian]{babel}
\usepackage{geometry}
\usepackage{indentfirst}
\usepackage{hyperref}
\usepackage{epigraph}
\usepackage{setspace}
\usepackage{parskip}
\usepackage{titlesec}
\usepackage{caption}

% Геометрия страницы (как в PDF: широкие поля, классический стиль)
\geometry{left=2.5cm,right=2cm,top=2cm,bottom=2cm}

% Настройка заголовков
\titleformat{\section}{\large\bfseries\centering}{}{0em}{}
\titleformat{\subsection}{\normalsize\bfseries}{}{0em}{}

% Отступы и межстрочный интервал
\setlength{\parindent}{1.25cm}
\setstretch{1.15}
\setlength{\parskip}{0.3\baselineskip}

% Убираем нумерацию страниц, но оставляем колонтитул по центру
\pagestyle{plain}

% Переопределение epigraph
\setlength{\epigraphwidth}{0.7\textwidth}
\renewcommand{\epigraphflush}{center}

\begin{document}

% Титул
\begin{center}
    {\Huge \textbf{ДАНИИЛ ХАРМС}} \\[1.5cm]
    {\huge \textbf{ВОПРОС}} \\[0.5cm]
\end{center}

\section*{ВОПРОС}

— Есть ли что-нибудь на земле, что имело бы значение и могло бы даже изменить ход событий не только на земле, но и в других мирах? — спросил я у своего учителя.

— Есть, — ответил мне мой учитель.

— Что же это? — спросил я.

— Это... — начал мой учитель и вдруг замолчал.

Я стоял и напряженно ждал его ответа. А он молчал.

И я стоял и молчал.

И он молчал.

И я стоял и молчал.

И он молчал.

Мы оба стоим и молчим.

Хо-ля-ля!

Мы оба стоим и молчим!

Хэ-лэ-лэ!

Да, да, мы оба стоим и молчим!

\begin{flushright}
\textit{16-17 июля 1937 г.}
\end{flushright}

\vspace{1cm}

\section*{ЗАБЫЛ, КАК НАЗЫВАЕТСЯ}

Один англичанин никак не мог вспомнить, как эта птица называется.

— Это, — говорит, — крюкица. Ах нет, не крюкица, а кирюкица. Или нет, не кирюкица, а курякица. Фу ты! Не курякица, а кукрикица. Да и не кукрикица, а кирикрюкица.

Хотите я вам расскажу рассказ про эту крюкицу? То есть не крюкицу, а кирюкицу. Или нет, не кирюкицу, а курякицу. Фу ты! Не курякицу, а кукрикицу. Да не кукрикицу, а кирикрюкицу! Нет, опять не так! Курикрятицу? Нет, не курикрятицу! Кирикрюкицу? Нет опять не так!

Забыл я, как эта птица называется. А уж если б не забыл, то рассказал бы вам рассказ про эту кирикуркукрукекицу.

\vspace{0.5cm}
\hrule
\vspace{0.5cm}

У одной маленькой девочки начал гнить молочный зуб. Решили эту девочку отвести к зубному врачу, чтобы он выдернул ей ее молочный зуб.

Вот однажды стояла эта маленькая девочка в редакции, стояла она около шкапа и была вся скрюченная.

Тогда одна редакторша спросила эту девочку, почему она стоит вся скрюченная, и девочка ответила, что она стоит так потому, что боится рвать свой молочный зуб, так как должно быть, будет очень больно. А редакторша спрашивает:

— Ты очень боишься, если тебя уколют булавкой в руку?

Девочка говорит:

— Нет.

Редакторша уколола девочку булавкой в руку и говорит, что рвать молочный зуб не больнее этого укола. Девочка поверила и вырвала свой нездоровый молочный зуб.

Можно только отметить находчивость этой редакторши.

\begin{flushright}
\textit{6 января 1937 г.}
\end{flushright}

\newpage

\section*{ПАШКВИЛЬ}

Знаменитый чтец Антон Исаакович Ш. — то самое историческое лицо, которое выступало в сентябре месяце 1940 года в Литейном лектории, — любило перед своими концертами полежать часок-другой и отдохнуть. Ляжет оно, бывало, на кушет и скажет:

— Буду спать, — а само не спит.

После концертов оно любило поужинать.

Вот оно придет домой, рассядется за столом и говорит своей жене:

— А ну, голубушка, состряпай-ка мне что-нибудь из лапши.

И пока жена его стряпает, оно сидит за столом и книгу читает.

Жена его хорошенькая, в кружевном передничке, с сумочкой в руках, а в сумочке кружевной платочек и ватрушечный медальончик лежат, жена его бегает по комнате, каблучками стучит, как бабочки, а оно скромно за столом сидит, ужина дожидается.

Все так складно и прилично. Жена ему что-нибудь приятное скажет, а оно головой кивает. А жена порх к буфетику и уже рюмочками там звенит.

— Налей-ка, душечка, мне рюмочку, — говорит оно.

— Смотри, голубчик, не спейся, — говорит ему жена.

— Авось, пупочка, не сопьюсь, — говорит оно, опрокидывая рюмочку в рот.

А жена грозит ему пальчиком, а сама боком через двери на кухню бежит.

Вот в таких приятных тонах весь ужин проходит. А потом они спать закладываются.

Ночью, если им мухи не мешают, они спят спокойно, потому что уж очень они люди хорошие!

\begin{flushright}
\textit{1940 г.}
\end{flushright}

\section*{УПАДАНИЕ}

Два человека упали с крыши пятиэтажного дома, новостройки. Кажется, школы. Они съехали по крыше в сидячем положении до самой кромки и тут начали падать.

Их падение раньше всех заметила Ида Марковна. Она стояла у окна в противоположном доме и сморкалась в стакан. И вдруг она увидела, что кто-то с крыши противоположного дома начинает падать. Вглядевшись, Ида Марковна увидела, что это начинают падать сразу целых двое. Совершенно растерявшись, Ида Марковна содрала с себя рубашку и начала этой рубашкой скорее протирать запотевшее оконное стекло, чтобы лучше разглядеть, кто там падает с крыши. Однако сообразив, что, пожалуй, падающие могут увидеть ее голой и невесть чего про нее подумать, Ида Марковна отскочила от окна за плетеный треножник, на котором стоял горшок с цветком.

В это время падающих с крыш увидела другая особа, живущая в том же доме, что и Ида Марковна, но только двумя этажами ниже. Особу эту тоже звали Ида Марковна. Она, как раз в это время, сидела с ногами на подоконнике и пришивала к своей туфле пуговку. Взглянув в окно, она увидела падающих с крыши. Ида Марковна взвизгнула и, вскочив с подоконника, начала спешно открывать окно, чтобы лучше увидеть, как падающие с крыши ударятся об землю. Но окно не открывалось. Ида Марковна вспомнила, что она забила окно снизу гвоздем, и кинулась к печке, в которой она хранила инструменты: четыре молотка, долото и клещи.

Схватив клещи, Ида Марковна опять подбежала к окну и выдернула гвоздь. Теперь окно легко распахнулось. Ида Марковна высунулась из окна и увидела, как падающие с крыши со свистом подлетали к земле.

На улице собралась уже небольшая толпа. Уже раздавались свистки, и к месту ожидаемого происшествия не спеша подходил маленького роста милиционер. Носатый дворник суетился, расталкивая людей и поясняя, что падающие с крыши могут вдарить собравшихся по головам.

К этому времени уже обе Иды Марковны, одна в платье, а другая голая, высунувшись в окно, визжали и били ногами.

И вот наконец, расставив руки и выпучив глаза, падающие с крыши ударились об землю.

Так и мы иногда, упадая с высот достигнутых, ударяемся об унылую клеть нашей будущности.

\begin{flushright}
\textit{7 сентября 1940 г.}
\end{flushright}

\newpage

\section*{ОЧЕНЬ СТРАШНАЯ ИСТОРИЯ}

\begin{verse}
Доедая с маслом булку, \\
Братья шли по переулку. \\
Вдруг на них из закоулка \\
Пес большой залаял гулко. \\

Сказал младший: «Вот напасть! \\
Хочет он на нас напасть. \\
Чтоб в беду нам не попасть, \\
Псу мы бросим булку в пасть». \\

Все окончилось прекрасно. \\
Братьям сразу стало ясно, \\
Что на каждую прогулку \\
Надо брать с собою… булку.
\end{verse}

\begin{flushright}
\textit{1938}
\end{flushright}

\end{document}